

\section{La Mathématique du Vocabulaire : 1100 lemmes = 80% de couverture}

L'enseignement du grec ancien, héritier d'une longue tradition humaniste, se trouve aujourd'hui
confronté à un paradoxe existentiel. D'un côté, la discipline conserve une aura de prestige
intellectuel et d'accès privilégié aux sources de la pensée occidentale ; de l'autre, elle fait face
à un constat d'échec pédagogique systémique : l'incapacité de la vaste majorité des apprenants, même
après plusieurs années d'études universitaires, à lire un texte original avec fluidité et sans
assistance constante. Les témoignages d'étudiants et d'enseignants convergent vers une réalité
décevante où le décodage laborieux, mot à mot, tient lieu de lecture, transformant l'expérience
esthétique et intellectuelle des textes en un exercice de cryptanalyse permanent.\cite{III-1-1-ref1}
Cette \og crise de la lecture\fg{} n'est pas imputable à un manque de rigueur ou de volonté, mais
elle signale une inadéquation structurelle entre les méthodes d'enseignement traditionnelles et la
réalité cognitive de l'acquisition d'une langue flexionnelle complexe.

C'est dans ce contexte que s'inscrit la section III.1.\cite{III-1-1-ref1} du plan de refondation,
intitulée \og La Mathématique du Vocabulaire\fg{}. Ce projet ne vise pas une simple réforme
marginale des pratiques, mais un changement de paradigme radical, fondé sur l'intégration des
données massives issues des humanités numériques et des acquis récents de la psycholinguistique. Le
postulat central est d'une simplicité trompeuse : la maîtrise du grec ancien est un problème
mathématique autant que philologique. L'analyse statistique des corpus, rendue possible par des
projets comme le \textit{Thesaurus Linguae Graecae} (TLG) et le \textit{Perseus Project}, a mis en
lumière une architecture lexicale précise, régie par des lois de distribution de fréquence
incontournables.

Ce rapport se propose d'explorer en profondeur la thèse selon laquelle un noyau restreint de lemmes
— environ 1 100 — permet de couvrir 80 \% de tout texte grec standard, et que la focalisation
exclusive sur ce noyau, au détriment des \textit{hapax legomena} (mots n'apparaissant qu'une fois)
et du vocabulaire rare, constitue la condition \textit{sine qua non} pour débloquer la lecture
fluide. Nous démontrerons que l'approche traditionnelle, qui tend à traiter tous les mots comme
ayant une valeur pédagogique égale, sature la mémoire de travail des étudiants et empêche
l'automatisation des processus de bas niveau nécessaires à la compréhension textuelle. En nous
appuyant sur la loi de Zipf, la théorie de la charge cognitive et l'analyse critique des manuels
actuels, nous validerons scientifiquement la nécessité de \og cesser d'apprendre des hapax\fg{} pour
se concentrer sur le \og Top 1000\fg{}, et nous esquisserons les contours d'une nouvelle pédagogie
de l'efficacité lexicale.



\subsection{Chapitre i : l'architecture statistique du corpus grec et la loi de zipf}

Pour refonder l'enseignement du vocabulaire, il est impératif de comprendre la \og physique\fg{}
interne de la langue grecque. Contrairement à une intuition pédagogique diffuse qui voudrait que
l'exposition aléatoire aux textes d'auteurs suffise à bâtir un vocabulaire compétent, la
linguistique de corpus révèle une inégalité structurelle et violente dans la distribution des mots.
Cette inégalité est régie par la loi de Zipf, un principe universel des langues naturelles qui
stipule que la fréquence d'un mot est inversement proportionnelle à son rang dans la table des
fréquences.\cite{III-1-1-ref3}



\subsection{La loi de zipf appliquée au grec ancien}

La loi de Zipf décrit une courbe de distribution hyperbolique : un très petit nombre de mots (les
mots de haute fréquence) accapare la majorité des occurrences textuelles, tandis qu'une immense \og
longue traîne\fg{} (\textit{long tail}) de mots rares ne contribue que marginalement à la couverture
du texte. En grec ancien, cette distribution prend une forme particulièrement accusée, favorable à
l'apprenant débutant mais redoutable pour l'expert.

Les analyses menées par Wilfred Major, s'appuyant sur la base de données morphologique du
\textit{Perseus Project} (plus de 4 millions de mots), ont permis de quantifier cette distribution
avec une précision inégalée. Il apparaît que le grec ancien possède un noyau lexical \og hyper-
dense\fg{}. Là où l'anglais ou le latin nécessitent un nombre plus élevé de lemmes pour atteindre
les premiers seuils de compréhension, le grec permet d'atteindre 50 \% de couverture textuelle avec
un nombre dérisoire de lemmes.

\begin{table}[h!]\small\centering\begin{tabular}{| l | l | l | l | l |}\hline\textbf{Seuil de Couverture} & \textbf{Grec Ancien (Lemmes)} & \textbf{Latin (Lemmes)} & \textbf{Anglais (Lemmes)} & \textbf{Implications Pédagogiques} \\ \hline\textbf{50 \%} & \~ 65 lemmes & \~ 200-250 lemmes & \~ 100 lemmes & Vocabulaire de survie immédiate (mots grammaticaux, verbes d'état). \\ \hline\textbf{80 \%} & \~ 1 100 lemmes & \~ 1 500 lemmes & \~ 2 300-3 000 lemmes & Seuil de compréhension globale de la structure (le \og Top 1000\fg{}). \\ \hline\textbf{90 \%} & \~ 3 020 lemmes & N/A & \~ 5 000 lemmes & Lecture confortable avec dictionnaire ponctuel. \\ \hline\textbf{95 \%} & \~ 5 948 lemmes & N/A & \~ 9 000 lemmes & Seuil d'autonomie réelle (Lecture extensive). \\ \hline\end{tabular}\end{table}Données synthétisées à partir de.\cite{III-1-1-ref5}

Ce tableau met en exergue l'avantage comparatif du grec ancien pour le débutant : avec seulement
\textbf{65 mots}, un lecteur identifie statistiquement \textbf{un mot sur deux} dans n'importe quel
texte.\cite{III-1-1-ref11} Ce \og noyau de 50 \%\fg{} est constitué quasi exclusivement de mots-
outils (articles, prépositions, conjonctions) et d'une poignée de verbes omnipotents (\textit{eimí},
\textit{gígnomai}, \textit{légō}, \textit{échō}).\cite{III-1-1-ref11} Cependant, la courbe s'aplatit
ensuite sévèrement. Pour passer de 80 \% à 95 \% de couverture (le seuil requis pour une lecture
sans dictionnaire selon Paul Nation), l'apprenant doit multiplier son stock lexical par cinq,
passant de 1 100 à près de 6 000 mots.\cite{III-1-1-ref9}



\subsection{La définition du noyau stratégique : les 1100 lemmes}

Le chiffre de \textbf{1 100 lemmes} n'est pas arbitraire. Il correspond au point d'inflexion optimal
où l'investissement mémoriel produit le rendement maximal en termes de couverture textuelle. Au-delà
de ce seuil, chaque nouveau mot appris apporte un gain de couverture marginal décroissant. Ce noyau,
identifié par Wilfred Major et corroboré par les travaux du \textit{Dickinson College Commentaries}
(DCC), constitue l'ossature syntaxique et sémantique de la langue.\cite{III-1-1-ref6}

La constitution de cette liste de 1 100 mots a nécessité des choix méthodologiques rigoureux pour
éviter les biais statistiques bruts. Wilfred Major a procédé à plusieurs ajustements sur les données
brutes de \textit{Perseus} :

1. \textbf{Exclusion des Noms Propres :} Les noms comme \textit{Athēnaîos} ou \textit{Sōkrátēs},
bien que très fréquents dans certains textes, ont été retirés pour garantir l'universalité de la
liste. Un apprenant ne doit pas encombrer sa mémoire à long terme avec des noms de personnages qui
peuvent être identifiés par leur majuscule.\cite{III-1-1-ref7}

2. \textbf{Traitement de la Polysémie et des Homographes :} Les outils automatiques peuvent
confondre des formes. Par exemple, des lemmes rares partageant une forme avec un mot très fréquent
(comme certains noms de vipères ressemblant à des formes verbales courantes) ont été épurés
manuellement.\cite{III-1-1-ref7}

3. \textbf{Réintégration Culturelle :} Environ 79 lemmes, statistiquement juste en dessous de la
barre des 80 \%, ont été réintégrés en raison de leur importance culturelle majeure (ex:
\textit{eleútheros} \- libre) ou de leur productivité dérivationnelle en anglais et
français.\cite{III-1-1-ref7}

Ce noyau de 1 100 mots se compose d'environ 404 verbes, 325 noms et 188 adjectifs, le reste étant
composé d'invariables.\cite{III-1-1-ref6} Il est crucial de noter que cette liste inclut les verbes
irréguliers les plus redoutables. En effet, la fréquence d'usage et l'irrégularité morphologique
sont fortement corrélées : les verbes les plus usités sont ceux qui ont résisté à l'analogie
régularisatrice au fil des siècles.\cite{III-1-1-ref14}



\subsection{La tyrannie des hapax legomena et la longue traîne}

Si le noyau de 1 100 mots représente la \og tête\fg{} de la distribution de Zipf, la \og longue
traîne\fg{} est constituée par les \textit{hapax legomena} (mots n'apparaissant qu'une seule fois
dans un corpus) et les \textit{dis legomena} (deux fois). C'est ici que se joue la bataille de
l'efficacité pédagogique.

Dans des corpus massifs, environ \textbf{40 \% à 60 \%} du vocabulaire distinct est constitué
d'hapax.\cite{III-1-1-ref3}

\begin{itemize}\item Dans l'Iliade d'Homère, environ un mot sur 18 est un hapax absolu.\cite{III-1-1-ref15}\item Dans le Nouveau Testament, près de 2 000 mots sur un total de 5 420 lemmes n'apparaissent qu'une seule fois.\cite{III-1-1-ref16}\item Plus globalement, une immense proportion des mots listés dans le dictionnaire \textit{Liddell-Scott-Jones} (LSJ) sont des mots que l'étudiant moyen ne rencontrera peut-être jamais, ou une seule fois dans sa vie.\end{itemize}La \og Mathématique du Vocabulaire\fg{} postule que l'apprentissage intentionnel de ces mots est une
erreur stratégique. Traditionnellement, l'enseignement philologique ne discrimine pas : lors de la
traduction d'un texte, l'étudiant cherche dans le dictionnaire un hapax avec la même diligence qu'un
mot fréquent, et tente souvent de le mémoriser. Or, le \og rendement\fg{} de cet effort est nul pour
les textes futurs. Passer 10 minutes à mémoriser un adjectif poétique rare qui ne reviendra jamais
est un gaspillage de temps cognitif (\textit{cognitive time on task}) qui aurait dû être alloué à la
consolidation d'un verbe du Top 500 dont la maîtrise est imparfaite.\cite{III-1-1-ref16}

La distribution des hapax n'est pas seulement une curiosité statistique ; c'est un obstacle majeur à
la fluidité. Comme le note James Tauber, même en connaissant les 500 mots les plus fréquents, on
laisse encore une part significative du texte dans l'ombre, une part composée majoritairement de ces
mots rares qui portent souvent la couleur locale et la précision sémantique du
passage.\cite{III-1-1-ref17} C'est pourquoi la stratégie ne doit pas être d'ignorer ces mots
\textit{dans le texte}, mais de cesser de les \textit{apprendre} (mémorisation active), pour les
traiter par des aides à la lecture (glossaires, outils numériques) qui n'encombrent pas la mémoire à
long terme.



\subsection{Chapitre ii : fondements cognitifs de la restriction lexicale}

La justification de la concentration sur le noyau de 1 100 lemmes dépasse la simple statistique ;
elle est enracinée dans la psychologie cognitive et les théories de l'acquisition du langage. Le
cerveau humain traite le langage selon des contraintes biologiques précises, notamment la capacité
limitée de la mémoire de travail et la nécessité d'automatiser les processus de bas niveau pour
libérer des ressources pour la compréhension de haut niveau.



\subsection{La théorie de l'efficacité verbale (verbal efficiency theory \- vet)}

La Théorie de l'Efficacité Verbale, formulée par Charles Perfetti, est centrale pour comprendre
pourquoi le vocabulaire fréquentiel est la clé de la lecture. La VET postule que la compréhension en
lecture est le produit de l'efficacité des processus d'identification des mots et d'analyse
syntaxique locale.\cite{III-1-1-ref19}

Le principe est celui d'un système à ressources limitées. La mémoire de travail (\textit{working
memory}) dispose d'une capacité finie. Si les processus de bas niveau (déchiffrage des lettres,
récupération du sens des mots, analyse morphologique) consomment une trop grande part de ces
ressources attentionnelles, il ne reste plus de capacité cognitive disponible pour les processus de
haut niveau : intégration sémantique, construction du sens global de la phrase, inférences logiques
et appréciation esthétique.\cite{III-1-1-ref22}

En grec ancien, le coût cognitif des processus de bas niveau est exorbitant pour un débutant :

1. \textbf{Décodage orthographique :} L'alphabet différent ralentit l'accès lexical.

2. \textbf{Ambiguïté morphologique :} Chaque mot doit être analysé (cas, genre, nombre, temps, mode)
avant d'être intégré syntaxiquement.

3. \textbf{Charge lexicale :} Si l'étudiant doit effectuer une recherche consciente en mémoire ou
dans un dictionnaire pour plus de 5 \% des mots, le processus de lecture s'effondre.

L'automatisation du noyau de 1 100 lemmes permet de contourner ce goulot d'étranglement. Lorsqu'un
mot est \og sur-appris\fg{} (\textit{overlearned}) et reconnu instantanément (automatisation), sa
récupération ne coûte presque rien à la mémoire de travail.\cite{III-1-1-ref24} Si 80 \% des mots
d'une phrase sont reconnus automatiquement, l'étudiant peut allouer toutes ses ressources
attentionnelles à la résolution des 20 \% restants (mots rares, hapax) et à la construction du sens
complexe de la phrase.\cite{III-1-1-ref26}



\subsection{L'hypothèse de la qualité lexicale (lexical quality hypothesis)}

L'Hypothèse de la Qualité Lexicale (LQH), également développée par Perfetti, affined la VET en se
concentrant sur la nature des représentations mentales des mots. La LQH suggère que la fluidité ne
dépend pas seulement de \textit{combien} de mots on connaît, mais de la \textit{qualité} de cette
connaissance.\cite{III-1-1-ref28}

Une représentation lexicale de haute qualité inclut une connexion forte et précise entre
l'orthographe, la phonologie et la sémantique du mot. À l'inverse, une représentation de \og basse
qualité\fg{} est floue, instable et lente à récupérer. Le problème de l'enseignement traditionnel
est qu'en forçant les étudiants à apprendre une masse indistincte de mots (y compris des hapax), on
crée des milliers de représentations de basse qualité. L'étudiant \og reconnaît vaguement\fg{}
beaucoup de mots mais doit souvent vérifier, hésite entre des sens proches, ou confond des
paronymes.

La stratégie du \og Top 1000\fg{} vise à construire des représentations de \textit{très haute
qualité} pour un nombre restreint de mots. En focalisant la répétition et l'exposition sur ces 1 100
lemmes, on s'assure que leur récupération est rapide, précise et résistante aux interférences, ce
qui est la définition même de la compétence lexicale selon la LQH.\cite{III-1-1-ref31}



\subsection{Théorie de la charge cognitive et "goulot d'étranglement"}

La théorie de la charge cognitive (Sweller) distingue trois types de charge : intrinsèque (la
difficulté inhérente de la tâche), extrinsèque (la difficulté liée à la manière dont l'information
est présentée) et germane (l'effort utile pour l'apprentissage).\cite{III-1-1-ref33}

L'apprentissage des hapax impose une \textbf{charge cognitive extrinsèque} massive. Présenter à un
étudiant des mots rares au même niveau visuel et pédagogique que des mots fréquents crée du \og
bruit\fg{}. Le cerveau de l'apprenant, ne pouvant distinguer l'utile du superflu, tente de tout
traiter avec la même intensité, ce qui mène rapidement à une surcharge (\textit{cognitive
overload}).\cite{III-1-1-ref35}

L'effet de goulot d'étranglement (bottleneck) est particulièrement visible dans le traitement
syntaxique. La mémoire de travail verbale doit maintenir les mots du début de la phrase en mémoire
tampon (buffer) jusqu'à ce que le verbe ou la fin de la clause permette l'intégration syntaxique
(surtout en grec où le verbe est souvent final ou l'hyperbate fréquente).\cite{III-1-1-ref37} Si la
récupération lexicale des mots intermédiaires est lente (non-automatisée), les mots du début de
phrase s'effacent de la mémoire tampon avant que la phrase ne soit comprise. C'est l'expérience
commune de l'étudiant qui, arrivé à la fin de la phrase grecque, a déjà oublié le
début.\cite{III-1-1-ref39}

L'automatisation du vocabulaire fréquentiel élargit fonctionnellement ce goulot d'étranglement en
accélérant le traitement, permettant à l'information de transiter plus vite que le taux d'effacement
de la mémoire à court terme.



\subsection{Le seuil critique de 95-98 \% pour l'input compréhensible}

Les recherches de Paul Nation et d'autres linguistes sur l'acquisition des langues secondes ont
établi des seuils précis pour la compréhension de texte.

\begin{itemize}\item \textbf{95 \% de couverture lexicale} est le seuil minimal pour une lecture assistée efficace, où le lecteur peut deviner le sens des mots inconnus grâce au contexte.\cite{III-1-1-ref10}\item \textbf{98 \% de couverture} est le seuil pour une lecture de plaisir, autonome et fluide.\cite{III-1-1-ref42}\end{itemize}Le noyau de 1 100 lemmes (80 \%) ne permet pas d'atteindre ces seuils immédiatement, mais il
constitue la plateforme de lancement indispensable. En dessous de 80 \%, le texte est un mur opaque.
À 80 \%, le texte devient un puzzle résoluble avec des aides. L'objectif de la refondation est
d'amener les étudiants le plus vite possible à ce plateau de 80 \%, pour ensuite utiliser des textes
adaptés (\textit{graded readers}) qui artificiellement augmentent ce pourcentage vers 95 \% en
limitant le vocabulaire rare.\cite{III-1-1-ref44}



\subsection{Chapitre iii : critique de la pédagogie traditionnelle au prisme des données}

L'adoption de la \og Mathématique du Vocabulaire\fg{} implique une remise en question profonde des
outils pédagogiques existants. L'analyse quantitative des manuels de grec ancien les plus populaires
révèle des déficiences majeures dans leur gestion du vocabulaire, déficiences qui expliquent en
grande partie l'échec de l'acquisition de la fluidité.



\subsection{Analyse comparative : athenaze et from alpha to omega}

L'étude séminale de Rachel Clark (2009), \og The 80\% Rule: Greek Vocabulary in Popular
Textbooks\fg{}, fournit des données accablantes sur l'alignement entre le vocabulaire enseigné et la
réalité statistique de la langue. Clark a comparé le vocabulaire de deux manuels standards,
\textit{Athenaze} et \textit{From Alpha to Omega}, avec les listes de fréquence de Wilfred
Major.\cite{III-1-1-ref46}

\begin{table}[h!]\small\centering\begin{tabular}{| p{2cm} | p{2.2cm} | p{2.2cm} | p{2.2cm} |}\hline\textbf{Manuel} & \textbf{Mots du \og Top 1100\fg{} introduits formellement} & \textbf{Pourcentage du Noyau Couvert} & \textbf{Mots introduits \+ Glossés dans les lectures} \\ \hline\textbf{From Alpha to Omega} & 463 lemmes & \textbf{42 \%} & 586 lemmes (53 \%) \\ \hline\textbf{Athenaze (Book I \& II)} & 602 lemmes & \textbf{54 \%} & 725 lemmes (66 \%) \\ \hline\end{tabular}\end{table}Source : Analyse des données de 46

Ces chiffres sont révélateurs. Un étudiant consciencieux qui termine \textit{From Alpha to Omega}
n'aura appris que \textbf{42 \%} des mots essentiels pour lire du grec réel. Même avec
\textit{Athenaze}, souvent considéré comme une méthode \og par la lecture\fg{}, un tiers du
vocabulaire critique reste inconnu à la fin du parcours.



\subsection{Le syndrome du vocabulaire "narratif" vs "fréquentiel"}

La raison de cette sous-performance réside dans la primauté donnée à la narration ou à la grammaire
sur la fréquence.

Dans Athenaze, l'histoire de Dikaiopolis le fermier impose un vocabulaire spécifique : \og
bœuf\fg{}, \og charrue\fg{}, \og champ\fg{}, \og navire\fg{}. Bien que charmants, ces mots ont une
fréquence globale faible dans le corpus de la littérature classique (Platon, Thucydide, Tragiques).
En revanche, des termes abstraits cruciaux, des verbes de pensée ou d'argumentation (nomizō,
hēgeomai), ou des conjonctions logiques complexes, moins utiles pour raconter une histoire de ferme
mais vitaux pour lire de la philosophie, sont sous-représentés ou introduits
tardivement.\cite{III-1-1-ref46}

De plus, Clark note un phénomène de \og Glossé et Oublié\fg{}. De nombreux mots à haute fréquence
apparaissent dans les textes de lecture des manuels mais sont simplement traduits en note de bas de
page (glossés) sans être inclus dans les listes de vocabulaire à mémoriser. L'étudiant les \og
voit\fg{} mais ne les \og apprend\fg{} pas, et faute de répétition systématique, ne les acquiert
jamais.\cite{III-1-1-ref46}



\subsection{Le problème de la séquenciation morphologique}

La méthode traditionnelle (Grammaire-Traduction) introduit le vocabulaire en fonction des catégories
grammaticales enseignées. On apprend les mots de la 1ère déclinaison, puis de la 2ème, etc. Cette
approche crée des aberrations statistiques.

Les verbes en \-mi (dídōmi, títhēmi, hístēmi, hiēmi) figurent tous dans le Top 50 des mots les plus
fréquents.\cite{III-1-1-ref11} Ils sont omniprésents. Pourtant, en raison de leur complexité
morphologique et de leur irrégularité, ils sont souvent relégués aux tout derniers chapitres des
manuels (chapitre 32 dans From Alpha to Omega). Cela signifie que l'étudiant passe des mois, voire
une année entière, sans maîtriser des verbes qu'il rencontrera dans presque chaque phrase de grec
réel. La logique grammaticale sabote ici la logique fréquentielle, empêchant l'automatisation
précoce du noyau dur de la langue.\cite{III-1-1-ref47}



\subsection{Chapitre iv : la stratégie du noyau – analyse et structure des 1100 lemmes}

La refondation pédagogique exige une connaissance intime de ce noyau de 1 100 lemmes. De quoi est-il
constitué? Comment se structure-t-il sémantiquement? C'est cette \og anatomie\fg{} du vocabulaire
fréquentiel qui doit guider la création de nouveaux curricula.



\subsection{Anatomie du "top 1100" : catégories et fonctions}

Le noyau de 1 100 lemmes se décompose en catégories grammaticales qui révèlent la nature
fonctionnelle de la langue. Contrairement aux listes thématiques (la maison, le corps humain), le
Top 1000 est dominé par les mots de structure et les verbes d'action générique.

\begin{itemize}\item \textbf{Verbes : \~ 404 lemmes (36 \%).} C'est la catégorie reine. La maîtrise de ces verbes est le défi principal. Il ne s'agit pas seulement de connaître le lemme (présent), mais les temps primitifs, car les verbes fréquents sont souvent polymorphes (supplétifs comme \textit{horaō} / \textit{eidon} / \textit{opsomai}). La liste des 50 \% contient à elle seule 8 verbes titanesques (\textit{eimí, légō, poiéō, échō}...).\cite{III-1-1-ref11}\item \textbf{Noms : \~ 325 lemmes (29 \%).} Équilibrés entre les déclinaisons, mais avec une forte prédominance de concepts abstraits (\textit{archē}, \textit{dynamis}, \textit{logos}) et sociaux (\textit{anthrōpos}, \textit{anēr}, \textit{basileus}) sur les objets concrets.\cite{III-1-1-ref14}\item \textbf{Adjectifs : \~ 188 lemmes (17 \%).} Incluant les quantificateurs essentiels (\textit{polys}, \textit{pas}, \textit{megas}).\cite{III-1-1-ref11}\item \textbf{Mots Fonctionnels (Invariables) : \~ 180-200 lemmes.} Bien que moins nombreux en type, ils représentent une part écrasante des \textit{tokens} (occurrences) dans le texte. Les particules (\textit{men}, \textit{de}, \textit{gar}, \textit{oun}), les prépositions et les pronoms constituent l'armature logique du grec. Leurs sens multiples et contextuels en font une priorité d'apprentissage absolue.\cite{III-1-1-ref12}\end{itemize}

\subsection{L'inflexion statistique : pourquoi s'arrêter à 1100?}

Le choix de 1 100 (ou 80 \%) n'est pas qu'une commodité ; il correspond à un seuil de rentabilité
(\textit{diminishing returns}).

\begin{itemize}\item \textbf{Niveau 1 (Top 65 \- 50\%) :} L'investissement est minime, le gain est immense. C'est la survie.\item \textbf{Niveau 2 (Top 500 \- \~65-70\%) :} C'est le niveau atteint par les bons manuels actuels. Suffisant pour des phrases simples, mais frustrant pour les textes littéraires.\cite{III-1-1-ref8}\item \textbf{Niveau 3 (Top 1100 \- 80\%) :} Le seuil critique. À ce stade, la structure de la phrase est claire. Les mots inconnus sont des \og îlots\fg{} entourés de contexte connu.\item \textbf{Niveau 4 (Top 3000 \- 90\%) :} Pour gagner ces 10 \% supplémentaires, il faut apprendre près de 2 000 mots de plus. L'effort est triplé pour un gain marginal de 10 \%.\end{itemize}La stratégie pédagogique rationnelle est donc de viser une maîtrise \textit{totale et automatisée}
du Top 1100. Au-delà, l'apprentissage doit changer de nature : il ne doit plus être intensif (par
cœur) mais extensif (rencontre au fil des lectures, sans obligation de rétention
immédiate).\cite{III-1-1-ref9}



\subsection{Variation dialectale et corpus : classique vs biblique}

Une nuance importante doit être apportée concernant la variation du corpus. Le \og Top 1100\fg{}
varie légèrement selon que l'on vise le grec classique (Attique) ou le grec biblique (Koinè du
Nouveau Testament).

\begin{itemize}\item Pour le \textbf{Nouveau Testament}, le noyau est encore plus concentré. Avec seulement \textbf{310 à 320 mots}, on atteint déjà près de \textbf{80 \%} de couverture, un chiffre remarquablement bas comparé au grec classique.\cite{III-1-1-ref16} Cela s'explique par la pauvreté lexicale relative et la répétitivité du texte biblique.\item Cependant, il existe un large recouvrement entre les deux listes (les mots fonctionnels et les verbes de base sont les mêmes). La refondation doit intégrer ces données : pour un étudiant en théologie, la cible prioritaire est le noyau biblique de 320 mots, puis l'extension vers le noyau classique de 1100 mots pour une culture plus large.\cite{III-1-1-ref51}\end{itemize}

\subsection{Chapitre v : vers une refondation pédagogique (stratégies et outils)}

La section III.1.\cite{III-1-1-ref1} ne se contente pas de poser un diagnostic ; elle appelle à une
action concrète. Comment enseigner ce noyau de 1 100 lemmes en évitant les écueils du passé? La
réponse réside dans l'alliance des méthodes communicatives, des lectures graduées et des outils
numériques.



\subsection{Stratégie de "triage" lexical et gestion des hapax}

L'enseignant doit devenir un gestionnaire de l'économie attentionnelle de l'étudiant.

1. \textbf{Exclusion Explicite :} Face à un texte, l'enseignant doit dire : \og Ce mot est un hapax.
Voici sa traduction. Ne l'apprenez pas.\fg{} Cette autorisation à l'oubli est libératrice. Elle
réduit la charge cognitive et focalise l'énergie sur ce qui compte.

2. \textbf{Focalisation Intensive :} À l'inverse, tout mot du Top 1100 rencontré doit être traité
avec rigueur. \og Ce verbe est dans le Top 500. Vous devez le connaître. Quelle est sa racine? Ses
temps primitifs?\fg{}

3. \textbf{L'outil Anki et la Répétition Espacée :} Pour garantir l'automatisation, l'usage de
logiciels de répétition espacée (\textit{Spaced Repetition Systems \- SRS}) comme Anki est
incontournable. Des paquets pré-construits basés sur les fréquences de Major ou du GNT permettent de
\og marteler\fg{} le noyau dur indépendamment des lectures.\cite{III-1-1-ref51}



\subsection{L'approche par les "graded readers" (lectures graduées)}

Puisque les textes authentiques sont trop difficiles (trop d'hapax), il faut fabriquer des textes
\og sas\fg{}. Les \textit{Graded Readers} sont des textes écrits en grec ancien mais limités
strictement à un niveau de vocabulaire donné.

\begin{itemize}\item Des initiatives modernes, comme les lecteurs de Christophe Rico ou les adaptations de classiques, créent des textes où 95 \% des mots appartiennent au Top 500 ou Top 1000.\cite{III-1-1-ref44}\item Ces textes permettent à l'étudiant de faire l'expérience de la lecture fluide \textit{avant} d'avoir le niveau pour lire Platon. Ils consolident le noyau par la rencontre en contexte (\textit{incidental learning}) et développent la vitesse de traitement.\cite{III-1-1-ref56}\end{itemize}

\subsection{Les méthodes communicatives : polis et biblingo}

La méthode de l'Institut Polis (Christophe Rico) et les plateformes comme Biblingo incarnent cette
refondation. Elles appliquent les principes de l'acquisition des langues vivantes au grec ancien.

\begin{itemize}\item \textbf{Polis Method :} En utilisant l'immersion orale et le TPR (\textit{Total Physical Response}), cette méthode favorise une internalisation profonde et intuitive du vocabulaire fréquentiel. Le vocabulaire est introduit non pas par déclinaison, mais par utilité pragmatique et fréquence (les verbes de mouvement, les impératifs courants). L'oralisation massive du noyau (parler, écouter) crée des traces mémorielles beaucoup plus robustes que la seule lecture visuelle.\cite{III-1-1-ref57}\item \textbf{Biblingo :} Cette plateforme numérique intègre explicitement les données de fréquence. Elle propose des textes et des exercices calibrés pour correspondre au niveau de vocabulaire acquis, assurant que l'étudiant reste toujours dans la zone de développement proximal (ZPD) où l'apprentissage est efficace.\cite{III-1-1-ref60}\end{itemize}

\subsection{Outils numériques et personnalisation du corpus}

Enfin, les humanités numériques offrent des outils de \og chirurgie\fg{} lexicale.

\begin{itemize}\item \textbf{The Bridge (Bret Mulligan) :} Cet outil permet de générer des listes de vocabulaire sur mesure pour un texte donné, en excluant les mots déjà connus ou les mots trop fréquents (pour se concentrer sur le vocabulaire spécifique d'un auteur) ou inversement. Il permet de visualiser exactement quels mots du noyau sont présents dans un texte cible.\cite{III-1-1-ref63}\item \textbf{Analyseurs Morphologiques (Alpheios, Perseus) :} Ils permettent une lecture \og just-in-time\fg{} où l'étudiant peut cliquer sur un hapax pour avoir le sens sans avoir à quitter le texte des yeux, minimisant la rupture de la charge cognitive.\cite{III-1-1-ref9}\end{itemize}

\subsection{Conclusion}

La section III.1.\cite{III-1-1-ref1} \og La Mathématique du Vocabulaire\fg{} constitue le pivot
d'une pédagogie du grec ancien enfin efficace. Elle ne propose pas de réduire la langue à une liste
de mots, mais de construire rationnellement les fondations sur lesquelles la maîtrise littéraire
pourra s'élever.

En acceptant la réalité statistique de la loi de Zipf et les contraintes cognitives de la mémoire de
travail, nous devons opérer un tri drastique :

1. \textbf{Sanctuariser le noyau de 1 100 lemmes} comme objectif prioritaire absolu, dont
l'automatisation conditionne tout le reste.

2. \textbf{Dévaloriser pédagogiquement les hapax} et le vocabulaire rare, en les traitant comme de
l'information transitoire assistée par des outils, et non comme du savoir à mémoriser.

3. \textbf{Réorganiser les cursus et les manuels} pour aligner la progression sur la fréquence et
non sur la morphologie traditionnelle.

C'est à ce prix que l'enseignement du grec ancien pourra sortir de l'impasse du déchiffrage
perpétuel et offrir à nouveau aux étudiants la joie de la lecture fluide, directe et vivante des
textes fondateurs.

